% Blurb on back cover, gets included in lulu cover as well
% as regular version, so be careful with formatting

{
\parindent=0pt
\parskip=\baselineskip
{\OPTbacktitlefont
\textit{摘自引言：}}
\OPTbackfont

\emph{同伦类型论}是一个新兴的数学分支，它以一种出人意料的方式，将若干看似不同的领域汇聚到了一起。这一学科基于近年来发现的\emph{同伦论}与\emph{类型论}之间的深刻联系。
它所触及的话题跨度之大，看似天差地别：从球面的同伦群，到类型检查的算法，再到弱 $\infty$-群胚的定义，无不在其范围之内。

同伦类型论还为数学的根基问题带来了新的理念。
一方面，我们有 Voevodsky （沃沃茨基）提出的精妙而优美的\emph{泛等公理（univalence axiom）}。
特别地，泛等公理蕴含了这样一条原则：同构的结构可以被等同，这正是数学家们在日常工作中一直乐于采用的原则，尽管它与传统基础理论的``官方''教条并不相容。
另一方面，我们还有\emph{高阶归纳类型（higher inductive types）}，它们以直接的、逻辑的方式描述了同伦论中一些基本的空间与构造：球面、柱面、截断、局部化等。
二者在经典的集合论基础中都无法直接刻画，但在同伦类型论中结合起来后，它们使一种全新的``同伦类型逻辑''成为可能。

由此浮现出一种新的数学基础的构想：它具有内在的同伦内容，以一种``在不变量意义下''理解数学对象的方式来组织数学；与此同时，它还可以方便地在机器上实现，从而为实际从事数学研究的人提供切实的助力。
这就是\emph{泛等基础（Univalent Foundations）}计划。

本书旨在对泛等基础的基本内容作首次系统性阐述，并汇集这种新型推理方式的若干实例——但既不要求读者了解或学习任何形式逻辑，也无需使用任何计算机证明助理。
我们相信，泛等基础终将成为集合论的一种可行替代方案，成为大多数数学家所从事的非形式化数学的``隐含基础''。

\bigskip

\begin{center}
  {\Large
  \textit{欲免费获取本书，请访问 HomotopyTypeTheory.org。}}
\end{center}
}
